% Halaman ini biasanya tidak memiliki nomor halaman atau menggunakan romawi
% \thispagestyle{empty} 

\begin{center}
    % --- JUDUL ---
    % Mengatur font size 14pt, bold, uppercase, dan spasi tunggal
    \begingroup
    \fontsize{14pt}{16.8pt}\selectfont % 14pt font size, 1.2x leading for single spacing
    \bfseries
    \MakeUppercase{Restorasi Dokumen Terdegradasi Menggunakan \\ Generative Adversarial Network dengan \\
    Diskriminator Dual-Modal dan Optimasi \\ Loss Function Berorientasi HTR}
    \par
    \endgroup

    \vspace{2.5cm} % Jarak vertikal

    % --- IDENTITAS MAHASISWA ---
    Oleh\\
    \vspace{0.3cm}
    \textbf{Jatniko Nur Mutaqin}\\
    \textbf{NIM: 23523314}\\
    \textbf{(Program Studi Magister Informatika)}

    \vspace{1.0cm}

    % --- INSTITUSI ---
    Institut Teknologi Bandung

    \vspace{2.5cm}

    % --- BAGIAN PERSETUJUAN ---
    Menyetujui\\
    Tim Pembimbing

    \vspace{0.5cm}
    
    % Titik-titik untuk tanggal
    Tanggal ..................................................

    \vspace{1.0cm}

    % --- TANDA TANGAN PEMBIMBING 1 (KETUA) ---
    Ketua

    \vspace{2.5cm} % Ruang untuk tanda tangan

    \underline{\hspace{7cm}} \\ % Garis bawah
    (Ir. I Gusti Bagus Baskara Nugraha, ST, MT, Ph.D)

    \vspace{1.5cm}

    % --- TANDA TANGAN PEMBIMBING 2 & 3 (ANGGOTA) ---
    % Menggunakan minipage untuk membagi layout menjadi dua kolom
    % \begin{minipage}[t]{0.45\textwidth}
    %     \centering
    %     Anggota
        
    %     \vspace{2.5cm} % Ruang tanda tangan
        
    %     \underline{\hspace{6cm}} \\
    %     (Nama Pembimbing 2)
    % \end{minipage}
    % \hfill % Mengisi ruang kosong di tengah agar minipage terdorong ke kiri & kanan
    % \begin{minipage}[t]{0.45\textwidth}
    %     \centering
    %     Anggota
        
    %     \vspace{2.5cm} % Ruang tanda tangan
        
    %     \underline{\hspace{6cm}} \\
    %     (Nama Pembimbing 3)
    % \end{minipage}

\end{center}